\begin{textsample}{NEG Dim 4 – 2009\_c4 – Score -6.00 – 2009\_c4/enade\_2009\_publi\_15f3034279.txt}  \label{ex:f4_neg_026}
Capítulo 4 Percepção sobre a Prova As percepções dos estudantes sobre a prova aplicada no ENADE/2009 na área de COMUNICAÇÃO SOCIAL – PUBLICIDADE E PROPAGANDA foram mensuradas por meio de nove itens que avaliaram desde o grau de dificuldade da prova até o tempo gasto para concluí -la.
As questões foram analisadas separando -se concluintes e \textbf{ingressantes}, e as percepções sobre a prova foram relacionadas com o desempenho dos estudantes e com a Região Geográfica de localização do curso6.
O desempenho dos estudantes dentro de cada Área/Subárea foi classificado em quartos ( 1/4 ) para os concluintes e para os \textbf{ingressantes} em separado ( ver Figura 1 para o esquema de representação gráfica destas grandezas ).
As notas finais foram colocadas em ordem ascendente e no primeiro quarto foram colocadas as 25 \% notas mais baixas, limitadas superiormente pelo percentil 25 que vem a ser o primeiro quartil.
No segundo quarto, as 25 \% seguintes, limitadas superiormente pelo percentil 50 que vem a ser a mediana.
No terceiro quarto, as 25 \% notas após as encontradas no segundo quarto, ou seja, as primeiras 25 \% acima da mediana, limitadas superiormente pelo percentil 75 que vem a ser o terceiro quartil.
No último quarto foram \textbf{alocadas} as 25 \% maiores notas.
Potencialmente, poderia existir uma relação entre a nota do estudante e a sua percepção sobre a prova.
Calculando as médias das respostas segundo os quartos de notas, seria possível reconhecer a existência ( ou não ) de algum \textbf{gradiente} nas respostas, por exemplo, uma melhor nota correspondendo a uma percepção menos negativa.
Este \textbf{gradiente} seria traduzido, por exemplo, no gráfico 4.1a, pela diminuição das áreas laranja e vermelha ( Difícil e Muito difícil ) com o aumento do desempenho ( indo do primeiro ao último quarto ).
A medida tradicional de “ grau de discriminação ” comparou os resultados dos estudantes do primeiro quarto ( 25 \% piores estudantes ) aos do último quarto ( 25 \% melhores ).
Esta comparação, no gráfico disponibilizado, foi idêntica à comparação da primeira barra ( primeiro quarto ) com a última barra ( último quarto ) de cada grupo 6 Curso, neste relatório, refere -se à unidade de análise para o Conceito ENADE e é caracterizado pela combinação de IES, Área/Subárea ( quando pertinente ) e município de habilitação. \textbf{ingressante} ou concluinte nos gráficos “ segundo quartos de nota final ” ( gráficos com números ímpares neste capítulo ).
Por outro lado, a disponibilização das informações segundo os quartos de nota permitiu verificar o comportamento de dois outros grupos intermediários que na análise somente dos extremos tem que ser suposto \textbf{monotônico}.
A seguir, serão apresentados os principais resultados relativos aos nove itens avaliados. 4.1 Grau de dificuldade da prova 4.1.1 Formação Geral Os Gráficos 4.1a, 4.1b e 4.1c apresentam a distribuição das respostas à questão “ Qual o grau de dificuldade desta prova na parte de Formação Geral? ” segundo quartos de nota final e grupo de estudantes concluintes e \textbf{ingressantes} que participaram do ENADE/2009 na área de COMUNICAÇÃO SOCIAL – PUBLICIDADE E PROPAGANDA.
No total de concluintes, 14,1 \% declararam que a prova na parte da Formação Geral foi Difícil/Muito difícil, e com um valor próximo para os \textbf{ingressantes}, 18,0 \%.
Tanto entre os concluintes, como entre os \textbf{ingressantes}, a proporção de indivíduos que classificaram esta parte como Difícil/Muito difícil diminui com o desempenho do estudante.
Por outro lado, a proporção que declarou a parte Fácil/Muito fácil foi maior entre os concluintes ( 23,5 \% ) do que entre os \textbf{ingressantes} ( 20,4 \% ).
Não se observou uma variação \textbf{monotônica} significativa desta proporção com o aumento no desempenho, nem entre os concluintes, nem entre os \textbf{ingressantes}.
Entre os concluintes, 62,4 \% classificaram esta parte com um nível médio de dificuldade.
A proporção entre os \textbf{ingressantes} foi ligeiramente menor, 61,6 \%.
Entre os concluintes e \textbf{ingressantes}, esta proporção apresentou um comportamento não \textbf{monotônico}, pois aumentou do 1º até o 3º quarto depois declina para o último quarto ; o maior percentual está no 3º quarto.
Os Gráficos 4.2a e 4.2b apresentaram a distribuição das respostas à questão “ Qual o grau de dificuldade desta prova na parte de Formação Geral? ” segundo a Região onde o curso está situado para os estudantes concluintes e \textbf{ingressantes} que participaram do ENADE/2009 na área de COMUNICAÇÃO SOCIAL – PUBLICIDADE E PROPAGANDA.
Todas as Regiões apresentaram estudantes.
Na Região Nordeste um percentual muito baixo ( 0,8 \% ) entre os estudantes concluintes declarou a prova como Muito difícil.
Entre os concluintes, a maior proporção de estudantes que considerou a parte de Formação Geral da prova como de dificuldade média entre todas as Regiões, foram os da Região Nordeste com 66,1 \%, já entre os \textbf{ingressantes}, foram os estudantes da Região Norte com 68,2 \%.
As Regiões com as menores proporções desta resposta foram, respectivamente, para os concluintes e \textbf{ingressantes}, Norte ( 61,4 \% ) e Sudeste ( 60,4 \% ). 4.1.2 Componente Específico Os Gráficos 4.3a, 4.3b e 4.3c apresentam a distribuição das respostas à questão “ Qual o grau de dificuldade desta prova na parte do Componente Específico? ” segundo os quartos de nota final e grupo de estudantes concluintes e \textbf{ingressantes} que participaram do ENADE/2009 na área de COMUNICAÇÃO SOCIAL – PUBLICIDADE E PROPAGANDA.
Em geral, para as desagregações consideradas, os estudantes classificaram o Componente Específico como mais Difícil que a Formação Geral.
No total de concluintes, 13,2 \% declararam que a prova na parte do Componente Específico foi Difícil/Muito difícil, contra um valor de quase o dobro para os \textbf{ingressantes}, 24,1 \%.
Tanto entre os concluintes como entre os \textbf{ingressantes}, a proporção de indivíduos que classificaram esta parte como Muito difícil diminui com o desempenho do estudante.
O mesmo comportamento foi notado para a proporção que declara Difícil/Muito difícil.
A proporção que declarou a parte Fácil/Muito fácil foi maior entre os concluintes ( 24,9 \% ) do que entre os \textbf{ingressantes} ( 15,1 \% ).
Entre os concluintes, conforme o esperado, a proporção dos que consideraram esta parte como Fácil cresce com o desempenho.
Entre os \textbf{ingressantes}, a proporção foi estável em torno de 12 \% independente do grau de desempenho na prova.
Entre os concluintes, 62,0 \% classificaram esta parte com um nível médio de dificuldade.
A proporção entre os \textbf{ingressantes} foi ligeiramente menor, 60,9 \%.
Entre os concluintes, esta proporção aumentou com o desempenho até o 3º quarto apresentando declínio para o último quarto.
Entre os \textbf{ingressantes}, esta proporção basicamente aumentou com o desempenho, indo de 55,0 \% no primeiro quarto de notas a 65,3 \% no último quarto.
Os Gráficos 4.4a e 4.4b apresentaram a distribuição das respostas à questão “ Qual o grau de dificuldade desta prova na parte do Componente Específico? ” segundo a Região onde o curso está situado para os estudantes concluintes e \textbf{ingressantes} que participaram do ENADE/2009 na área de COMUNICAÇÃO SOCIAL – PUBLICIDADE E PROPAGANDA.
Todas as Regiões apresentaram estudantes, tanto concluintes quanto \textbf{ingressantes}.
Entre os concluintes, a maior proporção de estudantes que considerou a parte do Componente Específico da prova como de dificuldade média entre todas as Regiões, foi na Região Nordeste com 65,2 \%.
Já entre os \textbf{ingressantes}, a maior proporção foi na Região Norte com 65,3 \%.
A Região com a menor proporção desta resposta para os concluintes foi a Região Sudeste ( 61,0 \% ) e para os \textbf{ingressantes} foram as Regiões Sul e Sudeste ( ambas com 60,4 \% ). 4.2 Extensão da prova em relação ao tempo total Os Gráficos 4.5a, 4.5b e 4.5c apresentam a distribuição das respostas à questão sobre a extensão da prova em relação ao tempo total, segundo quartos de nota final e grupo de estudantes concluintes e \textbf{ingressantes} que participaram do ENADE/2009 na área de COMUNICAÇÃO SOCIAL – PUBLICIDADE E PROPAGANDA.
Tanto para concluintes quanto para \textbf{ingressantes}, a maioria considerou a extensão da prova como adequada em relação ao tempo destinado à sua resolução, respectivamente 57,9 \% e 60,2 \%.
Entre os concluintes e \textbf{ingressantes} esta proporção cresceu com o desempenho.
A proporção dos que consideraram a extensão da prova Muito curta foi ligeiramente maior entre os concluintes ( 1,1 \% ) do que entre os \textbf{ingressantes} ( 1,0 \% ).
A proporção dos que consideraram a extensão da prova Curta/Muito curta foi maior entre os concluintes ( 5,5 \% ) do que entre os \textbf{ingressantes} ( 4,8 \% ).
Entre os concluintes, 0,5 \% dos estudantes com desempenho no 3º grupo declararam que a extensão da prova teria sido Muito curta com relação ao tempo total.
Entre os \textbf{ingressantes}, o menor percentual ocorre também no 3º quarto ( 0,3 \% ).
A maior proporção de concluintes declarando que a prova foi Curta/Muito curta foi entre os estudantes de pior desempenho, 7,2 \%.
A proporção dos que consideraram a extensão da prova Muito longa foi ligeiramente maior entre os concluintes ( 14,3 \% ) que entre os \textbf{ingressantes} ( 13,1 \% ).
A proporção dos que consideraram a extensão da prova Longa também foi ligeiramente maior entre os concluintes ( 22,4 \% ) do que entre os \textbf{ingressantes} ( 21,8 \% ).
Os Gráficos 4.6a e 4.6b apresentam a distribuição das respostas à questão sobre a extensão da prova em relação ao tempo total, segundo a Região onde o curso está situado para os estudantes concluintes e \textbf{ingressantes} que participaram do ENADE/2009 na área de COMUNICAÇÃO SOCIAL – PUBLICIDADE E PROPAGANDA.
Em todas as Regiões brasileiras, a maioria dos estudantes, tanto concluintes como \textbf{ingressantes}, perceberam a prova como Adequada.
Os concluintes da Região Sudeste ( 56,6 \% ) se situaram abaixo da média ( 57,9 \% ) para esta resposta, enquanto entre os \textbf{ingressantes}, os da Região Norte ( 65,1 \% ) e Sul ( 62,1 \% ) ficaram acima da média ( 60,2 \% ). 4.3 Compreensão dos Enunciados das Questões 4.3.1 Formação Geral Os Gráficos 4.7a, 4.7b e 4.7c apresentam a distribuição das respostas à questão “ Os enunciados das questões da prova na parte de Formação Geral estavam claros e objetivos? ” segundo quartos de nota final e grupo de estudantes concluintes e \textbf{ingressantes} que participaram do ENADE/2009 na área de COMUNICAÇÃO SOCIAL – PUBLICIDADE E PROPAGANDA.
Com relação aos enunciados das questões da parte de Formação Geral, as opiniões foram positivas : a grande maioria dos concluintes ( 72,3 \% ) e \textbf{ingressantes} ( 75,8 \% ) consideraram claros e objetivos os enunciados de todas ou da maioria das questões.
Entre os concluintes, esta proporção cresceu conforme o desempenho : 59,8 \% para o grupo com o pior desempenho a 79,7 \% para o de melhor.
Entre os \textbf{ingressantes} o mesmo comportamento se repete, esta proporção cresceu conforme o desempenho, indo de 66,1 \% a 85,0 \%.
Em linhas gerais, pode -se dizer que a combinação de Poucos e Nenhum \textbf{decresceu} com o desempenho tanto para concluintes como para \textbf{ingressantes} : cai de 20,7 \% a 7,0 \% para os concluintes e de 16,1 \% a 4,8 \% para os \textbf{ingressantes}.
Os Gráficos 4.8a e 4.8b apresentam a distribuição das respostas à questão “ Os enunciados das questões da prova na parte de Formação Geral estavam claros e objetivos? ” segundo a Região onde o curso está situado para os estudantes concluintes e \textbf{ingressantes} que participaram do ENADE/2009 na área de COMUNICAÇÃO SOCIAL – PUBLICIDADE E PROPAGANDA.
Em termos regionais, as proporções de \textbf{ingressantes} e de concluintes que consideraram claros e objetivos os enunciados de todas ou da maioria das questões não variaram muito ( no máximo 10 pontos percentuais ).
Entre os concluintes, os valores foram de 71,3 \% no Sudeste a 77,0 \% no Nordeste.
Entre os \textbf{ingressantes}, as Regiões extremas foram as mesmas e estas proporções variaram de 74,5 \% no Sudeste a 80,3 \% no Nordeste.
Com relação às respostas “ poucos ” e “ não, nenhum ” tanto para concluintes como para \textbf{ingressantes} as maiores proporções ocorreram na Região Sudeste, respectivamente 12,6 \% e 10,9 \%. 4.3.2 Componente Específico Os Gráficos 4.9a, 4.9b e 4.9c apresentam a distribuição das respostas à questão “ Os enunciados das questões da prova na parte de formação específica estavam claros e objetivos? ” segundo quartos de nota final e grupo de estudantes concluintes e \textbf{ingressantes} que participaram do ENADE/2009 na área de COMUNICAÇÃO SOCIAL – PUBLICIDADE E PROPAGANDA.
Com relação aos enunciados das questões da parte do Componente Específico, as opiniões foram também positivas : a grande maioria dos concluintes ( 75,6 \% ) e \textbf{ingressantes} ( 76,4 \% ) consideraram claros e objetivos os enunciados de todas ou da maioria das questões.
A proporção de concluintes declarando que todos os enunciados estavam claros e objetivos foi quase constante, em torno dos 22 \% para todos os quartos : com o percentual mais baixo sendo de 22,2 \% no último quarto e o mais elevado de 23,0 \% no 1º quarto.
Já a combinação de todas ou da maioria das questões foi crescente com o desempenho : indo de 61,6 \% para os com pior desempenho a 82,0 \% para o último quarto.
Entre os \textbf{ingressantes}, esta proporção também cresceu com o desempenho : indo de 67,7 \% para os com pior desempenho a 83,6 \% para o último quarto.
O número de estudantes que declararam que “ não, nenhum ” dos enunciados das questões da parte do Componente Específico era claro e objetivo, foi um pouco maior do que na parte de Formação Geral : 2,6 \% entre os concluintes e 2,0 \% entre os \textbf{ingressantes}.
Os que responderam “ poucos ” são 7,1 \%, tanto entre os concluintes como entre os \textbf{ingressantes}.
A proporção de estudantes que respondeu que cerca de metade dos enunciados da parte do Componente Específico era claro e objetivo \textbf{decresceu} com o desempenho, tanto para concluintes como para \textbf{ingressantes}. enunciados da parte do Componente Específico era claro e objetivo \textbf{decresceu} com o desempenho, tanto para concluintes como para \textbf{ingressantes}.
Os Gráficos 4.10a e 4.10b apresentam a distribuição das respostas à questão “ Os enunciados das questões da prova na parte de formação específica estavam claros e objetivos? ” segundo a Região onde o curso está situado para os estudantes concluintes e \textbf{ingressantes} que participaram do ENADE/2009 na área de COMUNICAÇÃO SOCIAL – PUBLICIDADE E PROPAGANDA.
Em termos regionais, as proporções de \textbf{ingressantes} e de concluintes que consideraram claros e objetivos os enunciados de todas ou da maioria das questões variaram menos entre os \textbf{ingressantes} ( no máximo 4 pontos percentuais ) do que entre os concluintes ( no máximo 5 pontos percentuais ).
Entre os concluintes, os valores foram de 74,6 \% no Sudeste a 79,3 \% no Nordeste.
Apenas a Região do Sudeste apresentou valor abaixo da média nacional ( 75,6 \% ) para os concluintes.
Com relação às respostas “ poucos ” e “ não, nenhum ” para concluintes a maior proporção ocorreu na Região Centro-Oeste com 10,8 \% - tendo a Região Sudeste valor ligeiramente abaixo ( 10,3 \% ) - e a menor proporção na Região Nordeste com 7,2 \%. 4.4 Suficiência das informações/instruções fornecidas Os Gráficos 4.11a, 4.11b e 4.11c apresentam a distribuição das respostas à questão “ As informações/instruções fornecidas para a resolução das questões foram suficientes para resolvê -las? ” segundo quartos de nota final e grupo dos estudantes concluintes e \textbf{ingressantes} que participaram do ENADE/2009 na área de COMUNICAÇÃO SOCIAL – PUBLICIDADE E PROPAGANDA.
Quando perguntados sobre as informações/instruções fornecidas para a resolução das questões, 84,1 \% dos \textbf{ingressantes} e 85,7 \% dos concluintes de todo o Brasil responderam que estas eram ou excessivas ou suficientes em todas ou na maioria das questões, demonstrando uma avaliação positiva da prova.
Levando -se em conta o desempenho dos participantes, sem grandes diferenças entre as opiniões de estudantes, \textbf{ingressantes} e concluintes, a proporção de estudantes com a opinião \textbf{supracitada} foi crescente com o desempenho, tanto para concluintes como para \textbf{ingressantes}.
Entre os concluintes a proporção vai de 76,9 \% para os de pior desempenho a 90,1 \% para os de melhor desempenho.
Entre os \textbf{ingressantes}, as proporções equivalentes são, respectivamente, 77,2 \% a 90,0 \%.
O complemento desta opinião, ou seja, os estudantes que responderam somente em algumas ou em nenhuma delas, apresentou uma proporção decrescente com o desempenho tanto para os concluintes como para os \textbf{ingressantes}.
A resposta “ não, em nenhuma ” delas apareceu em todos os quartos tanto para concluintes quanto para \textbf{ingressantes}, mas também com representatividade decrescente com o desempenho.
Os Gráficos 4.12a e 4.12b apresentam a distribuição das respostas à questão “ As informações/instruções fornecidas para a resolução das questões foram suficientes para resolvê -las? ” segundo a Região onde o curso está situado para os estudantes concluintes e \textbf{ingressantes} que participaram do ENADE/2009 na área de COMUNICAÇÃO SOCIAL – PUBLICIDADE E PROPAGANDA.
Os Gráficos 4.12a e 4.12b apresentam a distribuição das respostas à questão “ As informações/instruções fornecidas para a resolução das questões foram suficientes para resolvê -las? ” segundo a Região onde o curso está situado para os estudantes concluintes e \textbf{ingressantes} que participaram do ENADE/2009 na área de COMUNICAÇÃO SOCIAL – PUBLICIDADE E PROPAGANDA. 4.5 Dificuldade encontrada ao responder à prova Os Gráficos 4.13a, 4.13b e 4.13c apresentam a distribuição das respostas à questão “ Você se deparou com alguma dificuldade ao responder à prova?
Qual? ” segundo quartos de nota final e grupo de estudantes concluintes e \textbf{ingressantes} que participaram do ENADE/2009 na área de COMUNICAÇÃO SOCIAL – PUBLICIDADE E PROPAGANDA.
Perguntados sobre as dificuldades com as quais se depararam ao responder à prova, 29,5 \% dos \textbf{ingressantes} apontaram como principal dificuldade o desconhecimento do conteúdo ; valor um pouco mais alto da proporção dos que optaram pela forma diferente de abordagem do conteúdo ( 28,1 \% ).
Já os concluintes apontaram como principal dificuldade a forma diferente de abordagem do conteúdo, 32,1 \% ; seguido da falta de motivação para fazer a prova ( 30,8 \% ).
Na verdade para a quase totalidade das áreas, \textbf{ingressantes} apresentaram uma proporção de estudantes bem maior do que entre os concluintes declarando como principal dificuldade o desconhecimento da matéria.
Nesta área, o desconhecimento do conteúdo foi apontado somente por 9,6 \% dos concluintes, valores, grosso modo, decrescentes com o desempenho ( 12,6 \% no grupo de pior desempenho e 8,8 \% no de melhor ).
Entre os concluintes, o comportamento foi inverso, a proporção de estudantes que apontaram o desconhecimento do conteúdo foi decrescente com o desempenho.
A falta de motivação foi apontada por 21,5 \% dos \textbf{ingressantes}.
Esta dificuldade apresentou valores decrescentes entre os concluintes caindo de 34,1 \% no primeiro quarto para 28,5 \% no último quarto.
Entre os \textbf{ingressantes}, os valores foram igualmente decrescentes : o grupo de pior desempenho apresentou o maior percentual com 25,1 \%.
Os Gráficos 4.14a e 4.14b apresentam a distribuição das respostas à questão “ Você se deparou com alguma dificuldade ao responder à prova?
Qual? ” segundo a Região onde o curso está situado para os estudantes concluintes e \textbf{ingressantes} que participaram do ENADE/2009 na área de COMUNICAÇÃO SOCIAL – PUBLICIDADE E PROPAGANDA.
Em relação às Regiões brasileiras, o percentual de \textbf{ingressantes} que apontaram o desconhecimento do conteúdo como principal dificuldade ao responder à prova variou de 27,9 \% na Região Centro-Oeste a 35,4 \% na Região Sul.
Entre os concluintes, tal percentual variou de 9,1 \% na Região Nordeste a 12,7 \% na Norte.
A forma diferente de abordagem de conteúdo foi citada como maior dificuldade por 24,2 \% dos \textbf{ingressantes} na Região Nordeste chegando a 34,7 \% na Região Norte.
Já entre os concluintes, os valores foram de 28,6 \% na Região Nordeste a 37,5 \% na Região Norte.
O percentual de concluintes que citaram a Falta de Motivação como principal dificuldade variou de 23,4 \% na Região Norte a 33,4 \% na Região Nordeste.
Entre os \textbf{ingressantes}, as estatísticas correspondentes foram de 17,8 \% na Região Norte a 24,3 \%, na Região Nordeste.
Os que declararam que não tiveram qualquer tipo de dificuldade variaram, entre os concluintes, de 17,2 \% na Região Centro-Oeste a 23,1 \% na Região Sul.
Entre os \textbf{ingressantes}, este percentual variou de 13,2 \% na Região Sul a 17,3 \% na Região Nordeste. 4.6 Percepção sobre os conteúdos das questões objetivas da prova Os Gráficos 4.15a, 4.16b e 4.16c apresentam a distribuição das respostas à questão “ Considerando apenas as questões objetivas da prova, você percebeu que :... ” segundo quartos de nota final e grupo de estudantes concluintes e \textbf{ingressantes} que participaram do ENADE/2009 na área de COMUNICAÇÃO SOCIAL – PUBLICIDADE E PROPAGANDA,.
As respostas possíveis para esta questão estavam relacionadas à cobertura do conteúdo das questões objetivas durante o curso.
Como esperado, a proporção de \textbf{ingressantes} ( 41,9 \% ) que declarou que não estudou ainda a maioria desses conteúdos foi bem maior do que entre os concluintes ( 6,8 \% ).
Esta proporção não apresentou comportamento \textbf{monotônico} como função do desempenho, nem para os concluintes nem para os \textbf{ingressantes}.
Por outro lado, a proporção de estudantes que declarou que estudou e aprendeu todos ou muitos desses conteúdos, entre os \textbf{ingressantes} foi de 38,6 \% e entre os concluintes 77,9 \%.
Esta proporção foi crescente como função do desempenho, tanto para os concluintes como para os \textbf{ingressantes}.
A proporção de estudantes que declararam ter estudado a maioria desses conteúdos, mas não os ter aprendido é maior entre os concluintes ( 9,8 \% ) do que os \textbf{ingressantes} ( 6,6 \% ). 4.6 Percepção sobre os conteúdos das questões objetivas da prova Os Gráficos 4.15a, 4.16b e 4.16c apresentam a distribuição das respostas à questão “ Considerando apenas as questões objetivas da prova, você percebeu que :... ” segundo quartos de nota final e grupo de estudantes concluintes e \textbf{ingressantes} que participaram do ENADE/2009 na área de COMUNICAÇÃO SOCIAL – PUBLICIDADE E PROPAGANDA,.
As respostas possíveis para esta questão estavam relacionadas à cobertura do conteúdo das questões objetivas durante o curso.
Como esperado, a proporção de \textbf{ingressantes} ( 41,9 \% ) que declarou que não estudou ainda a maioria desses conteúdos foi bem maior do que entre os concluintes ( 6,8 \% ).
Esta proporção não apresentou comportamento \textbf{monotônico} como função do desempenho, nem para os concluintes nem para os \textbf{ingressantes}.
Por outro lado, a proporção de estudantes que declarou que estudou e aprendeu todos ou muitos desses conteúdos, entre os \textbf{ingressantes} foi de 38,6 \% e entre os concluintes 77,9 \%.
Esta proporção foi crescente como função do desempenho, tanto para os concluintes como para os \textbf{ingressantes}.
A proporção de estudantes que declararam ter estudado a maioria desses conteúdos, mas não os ter aprendido é maior entre os concluintes ( 9,8 \% ) do que os \textbf{ingressantes} ( 6,6 \% ).
Os Gráficos 4.16a e 4.16b apresentam a distribuição das respostas à questão “ Considerando apenas as questões objetivas da prova, você percebeu que :... ” segundo a Região de habilitação do curso para os estudantes concluintes e \textbf{ingressantes} que participaram do ENADE/2009 na área de COMUNICAÇÃO SOCIAL – PUBLICIDADE E PROPAGANDA, segundo a região onde o curso esta situado.
Em relação às Regiões, concluintes que afirmaram não ter estudado ainda a maioria dos conteúdos ocorreram em todas as Regiões.
Entre os \textbf{ingressantes}, os percentuais variaram de 39,5 \% na Região Sudeste a 48,6 \%, na Região Norte.
Em todas as Regiões, a maioria dos concluintes afirmaram ter estudado e aprendido todos ou muitos dos conteúdos, com percentuais que variaram de 73,4 \% na Região Centro-Oeste a 82,2 \%, na Nordeste.
Para os \textbf{ingressantes}, tal percentual variou de 34,9 \%, na Região Sul, a 39,8 \%, na Centro-Oeste. 4.7 Tempo gasto para concluir a prova Os Gráficos 4.17a, 4.17b e 4.17c apresentam a distribuição das respostas à questão “ Qual foi o tempo gasto por você para concluir a prova? ” segundo quartos de nota final e grupo de estudantes concluintes e \textbf{ingressantes} que participaram do ENADE/2009 na área de COMUNICAÇÃO SOCIAL – PUBLICIDADE E PROPAGANDA.
Quando perguntados quanto ao tempo de conclusão da prova, a maioria respondeu “ Entre uma e duas horas ”, tanto entre os concluintes ( 48,3 \% ) quanto entre os \textbf{ingressantes} ( 48,2 \% ).
A segunda maior incidência de respostas foi “ Entre duas e três horas ”, igualmente para ambos os grupos : concluintes ( 29,5 \% ) e \textbf{ingressantes} ( 30,7 \% ).
A proporção daqueles que declararam ter usado às quatro horas, mas não ter conseguido terminar foi bem baixa : 1,3 \% entre os \textbf{ingressantes} e entre os concluintes.
Esta resposta ficou concentrada quase que exclusivamente entre os estudantes do primeiro quarto, aquele com pior desempenho.
Tanto entre os concluintes, como entre os \textbf{ingressantes} percebeu -se uma tendência crescente do tempo necessário para concluir a prova como função do desempenho.
Os Gráficos 4.18a e 4.18b apresentam a distribuição das respostas à questão “ Qual foi o tempo gasto por você para concluir a prova? ” segundo a Região onde o curso está situado para os estudantes concluintes e \textbf{ingressantes} que participaram do ENADE/2009 na área de COMUNICAÇÃO SOCIAL – PUBLICIDADE E PROPAGANDA.
A proporção de concluintes que gastaram “ Entre uma e duas horas ” foi maior na Região Sul ( 52,6 \% ) e menor na Região Centro-Oeste ( 45,0 \% ).
Para \textbf{ingressantes}, o maior percentual para este tempo ocorreu na Região Sul ( 52,0 \% ) e menor na Região Norte ( 43,5 \% ).
Considerando -se a distribuição cumulativa, entre os concluintes, os estudantes da Região Nordeste foram os que levaram mais tempo para concluir a prova.

% matched lemmas: alocar, decrescer, gradiente, ingressante, monotônico, supracitado
\end{textsample}
